\documentclass[12pt,fontset=fandol]{ctexart}

\usepackage[a4paper,margin=2.5cm]{geometry}
\usepackage{amsmath,amssymb}
\usepackage{graphicx}
\usepackage{hyperref}
\usepackage{xcolor}

\hypersetup{colorlinks=true, linkcolor=blue!50!black, urlcolor=blue!50!black}

\title{中文排版示例}
\author{你的名字 \\ 某某大学}
\date{\today}

\begin{document}

\maketitle

\begin{abstract}
  这是一段中文摘要示例。中文文档需要使用 \texttt{ctex} 宏集（或 \texttt{xeCJK}）并配合中文字体，
  否则编译时会提示 “Missing character”。
\end{abstract}

\section{引言}

中文段落之间的空行用于分段。行内公式同样可用：$a^2 + b^2 = c^2$。

\section{公式与图表}

\begin{equation}
  f(x) = \frac{1}{\sqrt{2\pi}\sigma}\exp\left(-\frac{(x-\mu)^2}{2\sigma^2}\right).
\end{equation}

\section{结论}

到这里，一份中文文档就完成了。

\end{document}
